% !TeX encoding = UTF-8 Unicode
% vim: set fileencoding=utf-8 :
\documentclass[11pt,twoside,openany]{scrbook}
\usepackage[a5paper,includeheadfoot,inner=18mm,outer=14mm,top=16mm,bottom=22mm,bindingoffset=6mm]{geometry}
\usepackage[T1]{fontenc}
\usepackage[utf8]{inputenc}
\usepackage[ngerman]{babel}
\usepackage{mathpazo}
\linespread{1.20}
\frenchspacing
\setlength{\parindent}{1.25em}
\setlength{\parskip}{0pt}
\clubpenalty=10000
\widowpenalty=10000
\displaywidowpenalty=10000

% Diese Datei ist UTF-8-kodiert.
\newcommand{\romantitel}{Die Raumfalte}
\newcommand{\romanautor}{Goetz Markgraf}

%%% BEGIN DOCUMENT
\begin{document}
%\nonfrenchspacing
%\sffamily

\begin{titlepage}
  \thispagestyle{empty}
  \centering
  \vspace*{0.22\textheight}

  \vspace{2.5em}
  {\huge\bfseries \romantitel\par}
  \vspace{3em}
  {\Large \romanautor\par}

  \vfill
\end{titlepage}

\tableofcontents

\chapter*{Ideen}

Es gelangen insgesamt 21 vermutlich automatisierte Riesenraumsonden in die Galaxis. Der Ursprung ist unbekannt, der Zweck ebenso. Zuerst wird die Signatur eines Schiffes bemerkt. Tamara Svårregad (Tochter des Piloten) findet ebenfalls die Signatur und vermutet, sie käme aus Richtung Virgo Haufen. Diese Hypothese im Rücken findet sie weitere fünf Signaturen. Später werden aber noch weitere Signaturen und auch Schiffe gefunden, die offenbar aus anderen Richtungen gekommen sind.




\chapter{Aufstieg der Monarchie}

Letztlich war es der Wohlstand, der das Ende der demokratischen Periode einläutete. Je besser es allen Menschen in ihrem persönlichen Umfeld ging, desto weniger haben sie sich dafür interessiert, was in der großen Politik geschah, so lange ihr persönlicher Wohlstand nicht angetastet wurde. ">Geiz ist geil!"< lautete damals ein sehr erfolgreicher Werbespruch, der eine Menge über die vorherrschende Mentalität aussagte.

Geiz, Wohlstand und Reichtum waren wichtiger geworden als Selbstbestimmung in der Politik. Das ging damals auch einher mit einer Angleichung der politischen Lager. Das Volk interessierte sich nicht mehr so sehr für Politik, weil diese auch mehr und mehr austauschbar geworden war. Ob rechts oder links, ob liberal oder gemäßigt, alle großen Parteien taten nach der Wahl mehr oder weniger das gleiche. ">Politikverdrossenheit"< war damals der Fachbegriff dafür. Die Wahlbeteiligung ging mehr und mehr zurück, der Wahlkampf wurde permanent aufwändiger, es wurden Unterschiede und Abgrenzungen zwischen den Parteien herausgestellt, die schon längst nicht mehr existieren. Die Demokratie war im Niedergang begriffen.

Wie so häufig, in solchen Zeiten gibt es jemanden, der die Zeichen der Zeit versteht und sie nutzen kann. Dieser Jemand war in diesem Fall General Padisha von Monroe. Der General war in mehreren Armeen ausgebildet worden. Aus welchem Land er ursprünglich entstammte, lies sich später nicht mehr klären, denn nach der Machtergreifung hatte in beispielloser Art und Weise einen Personenkult um seine Person und seine Familie errichtet, der die historischen Tatsachen verschleierte. Man kann aber als gesichert annehmen, dass er aus einer europäischen Familie stammte und dort eine umfangreiche Ausbildung genossen hatte. Dann hatte er vermutlich in mehreren afrikanischen Ländern in Söldnerarmeen gedient und diese befehligt. Ebenfalls kann vermutet werden, dass er einige Zeit in fundamentalistisch--islamischen Ländern verbracht hat und die dortige Mentalität aufs genaueste studierte.

Als dann die Beteiligungen bei Parlamentswahlen von gleich fünf europäischen Ländern bei unter 30\,\% lag, das Europaparlament aufgrund zu geringer Beteiligung sich als nicht handlungsfähig erklärte und der amerikanische Präsident nur mit viermaliger Stimmenauszählung ermittelt werden konnte, wähnte von Monroe seine Zeit als gekommen. Ohne selber über eine große Armee zu gebieten, übernahm er in einem Putsch die europäische Zentralregierung und setze sich als erster Befehlshaber des vereinten Kontinentes ein. Die jeweiligen Präsidenten, Kanzler und Premier Minister lies er zunächst im Amt, gab ihnen aber sehr konkrete Anweisungen. Mit dem amerikanischen Präsident schloss er zunächst nur einen Vertrag. Erst eine Wahlperiode später wurde er auch dort zum obersten Regierungschef gewählt. Der russische Präsident übergab sehr schnell die Machtbefugnisse im Land an den General und man munkelt, es wären größere Mengen Alkohol im Spiel gewesen. Da er nun über eine deutlich größere Armee verfügte, putschte er in fast allen afrikanischen Staaten, teilweise in Person, teilweise über Strohmänner, denen er befahl.

Danach arrangierte er sich mit den großen Scheich-Familien des vorderen Orients, gab ihnen spezielle Machtbefugnisse und Einflussbereiche, übernahm aber selber vollständig die Kontrolle der jeweiligen Politik des Landes. Und so ging es weiter. Noch bevor fünf Jahre vorbei waren, war General von Monroe entweder direkt oder indirekter Regierungschef von 85\,\% der Weltbevölkerung. Die größte Hürde hatten die Chinesen dargestellt, die grundsätzlich eine nicht chinesische Regierung abgelehnt hatten. Hier hatte die mittlerweile deutliche militärische Übermacht des Generals die Obersten des Landes dazu gebracht, die Selbstständigkeit aufzugeben. Immerhin gab es auch dafür Präzedenzfälle in der chinesischen Geschichte: die Yuan-Dynastie war schließlich auch nicht chinesisch, sondern mongolisch gewesen und der berühmte letzte Kaiser Pu~Yi seinerseits Manschure. Fremdherrschaft war also auch in China etwas, was man schon einmal erlebt hatte.

Die darauf folgende ">Aufräumarbeit"< dauerte noch deutlich länger, aber schon zu diesem Zeitpunkt war der Machtanspruch von General von Monroe nicht mehr gefährdet.

Und was tat das Volk in der Zeit? Gab es Aufstände, Widerstände, Rebellion? Nichts dergleichen! In den europäischen Ländern, in Nordamerika und auch in Japan war es der Bevölkerung völlig egal. Wichtig war doch nur, dass jeder einzelne weiterhin seinen persönlichen Wohlstand hatte und darauf achtete von Monroe sehr genau. Den meisten Bürgern ging es schon kurze Zeit nach der Übernahme deutlich besser als vorher, denn von Monroe setzte überall eine deutlich effizientere Verwaltung ein und Unmengen von Steuergeldern, die in den bürokratischen Apparaten versickert waren, konnten nun gewinnbringender eingesetzt werden. Durch geschickte Finanz- und Währungspolitik wurden die meisten dieser Länder bald schuldenfrei.

Ähnlich erging es auch den islamischen, südamerikanischen und afrikanischen Ländern. Hier hatte vorher in großem Stil die Korruption regiert. Unter Anwendung härtester Strafen verfolgte der General korrupte Beamte und entfernte sie aus dem Apparat. Statt dessen wurde massiv das Allgemeinwohl der Bevölkerung gesteigert, um der vorherrschenden fanatischen Neigung kleiner aber einflussreicher Bevölkerungsteile den Nährboden zu nehmen. Natürlich ging das nicht ohne Widerstand durch die zuvor herrschenden Familien und Clans. Da diese aber bereits entmachtet waren und über keinerlei Rückhalt in der Bevölkerung mehr verfügten, konnten sie im Keim erstickt werden.

Dennoch beäugte gerade die intellektuelle Oberschicht jedes aufgeklärten Landes die neue Regierung, die schnell die Züge einer Alleinherrschaft zeigte. Doch behielt von Monroe die bislang erfolgreiche Gewaltenteilung bei, er wurde nur der Vorsitzende der Legislative, die Exekutive und Jurisprudenz unterstand ihm nicht direkt. Damit waren weite Teile der argwöhnischen Gebildeten beruhigt.

Ein anderer Kunstgriff von Monroes betraf die Religion. Er hatte erkannt, dass sie immer noch für breite Teile der Weltbevölkerung wichtig ist. Daher stellte er sich in eine vage, nicht näher ausformulierte Tradition der jüdisch--christlich--islamischen Religion, indem er durchaus den einen Gott anerkannte und sich später sogar als ">von Gottes Gnaden"< titulierte. Andere Religionen lies aber weiterhin zu und arrangierte es so geschickt, dass er in allen diesen Religionen als Mittler zwischen dem Menschen und der himmlischen Sphäre eingesetzt war. So wurde er Regierungschef und oberster Priester der ganzen Weltbevölkerung.

Niemand kann einen Krieg oder eine andere Änderung der bestehenden Verhältnisse gegen den Willen der Industrie durchführen. Das wusste auch von Monroe. So kam er den Großkonzernen gar nicht erst in die Quere. Im Gegenteil: die Industrie war immer daran interessiert, dass es gut betuchte Konsumenten gibt, die die angebotenen Produkte kauften. Und durch den allgemeinen Wohlstand, der von Monroe schon alleine deshalb am Herzen lag, weil es die Massen ruhig hielt, wurde auch dieser Wunsch erfüllt. Außerdem sind Konzerne immer an stabilen Verhältnissen interessiert und durch die neue Regierungsform zeichnete sich ab, dass nicht alle drei bis fünf Jahre eine neue Regierung an die Macht kommen würde, die mit neuen Gesetzen und Vorschriften die Arbeit der Firmen erschweren würde. So unterstützten die meisten Konzerne den General -- wenn auch nicht offen, so aber durchaus im Verborgenen.

Manche hatten erwartet, eine Militärdiktatur erwachsen zu sehen, doch von Monroe hatte andere Pläne und so rief er sich schließlich im zehnten Jahr nach der Machtergreifung in Europa und Nordamerika zum \emph{Kaiser Padisha I.} des terranischen Reiches aus. Die jeweiligen Landesregierungen wurden nach und nach in eine neue Adelsschicht des Reiches umgewandelt. Die Monarchie war zurückgekehrt, allerdings in neuer Form, der Form der sog. ">aufgeklärten Monarchie"<. Damit hatte niemand gerechnet.

Doch die Analyse von Monroes war absolut richtig gewesen. Eine Diktatur würde zu viele Erinnerungen an Nazideutschland, den Duce aus Italien und andere derartige Herrschaften wecken. Die jeweiligen vergangenen Monarchen waren aber fast überall noch in guter Erinnerung. Selbst Napoleon, der sicherlich genug Greuel verübt hatte, um die obige Reihe aufgenommen zu werden, war den Franzosen in bester Erinnerung. Sicherlich, weil er den Einfluss der \emph{Grand Nation} zu neuem Ruhm verholfen hat. Aber vielleicht auch, weil er sich schließlich als König der Franzosen gekrönt hatte.

Seine Abstammung behielt Padisha von Monroe sehr für sich. Er sprach ein Dutzend Sprachen fließend. Seine äußere Erscheinung lies sich nicht eindeutig einem Phänotyp zuordnen, er konnte ein dunkelhäutiger Europäer oder ein hellhäutiger Nordafrikaner sein. Seine ausgeprägten Backenknochen machten ihn sogar -- wenn er es darauf anlegte -- fast zu einem Indio Südamerikas und der leicht struppige schwarze Bartwuchs lies ihn leicht arabisch erscheinen. Nur als Chinese oder Japaner ging er nicht durch, doch machte er dieses Manko durch perfekten Gebrauch der jeweiligen Landessprachen wieder wett. Es gab sogar das Gerücht, dass von Monroe nicht etwa natürlich gezeugt wäre, sondern als genmanipulierter Mensch das Erbgut der ganzen Menschheit in sich träge. Vermutlich wurde dieses Gerücht durch die extrem gut funktionierenden Propagandaabteilungen in die Welt gesetzt, es war allerdings auch nicht völlig ausgeschlossen, dass es der Wahrheit entsprach.

Wie auch immer, jede Nation konnte zum General aufschauen und sich selbst darin erblicken. Auf diese Weise hatte jede Nation das Gefühl, selber den Kaiser gestellt zu haben. Zweifelsohne war das ein Effekt, den von Monroe beabsichtigt hatte oder zumindest auszunutzen verstand.

So waren die politischen Verhältnisse geklärt: An der Spitze der Regierung stand ein sog. ">wohlwollender Alleinherrscher"<, der das Volk nicht unterdrückte und auch nicht ausbeutete, gleichwohl es ihm sicher nicht schlecht ging. Die wirtschaftliche Lage war langfristig geklärt, dem Volk ging es gut, meist besser als vor der Übernahme und ein Widerstand war nicht in Sicht. Die umfangreichste Machtübernahme der Welt war mit einem Mindestmaß an Gewalt abgelaufen; kaum ein Schuss war gefallen und insgesamt waren in kleineren Scharmützeln in entlegenen Teilen kaum mehr als 1.000 Menschen ums Leben gekommen. Der ">sanfte Putsch"< wurde das später hinter vorgehaltener Hand genannt, mehr eine Schmeichelei als eine Verwünschung.

Neben den vielen Veränderungen, etwa die der Straffung der Regierungsapparate, den Umbau der Demokratien zu neuen Feudalherrschaften oder dem Einsetzen der allgemeinen Wehrpflicht \underline{und} der zivilen Arbeitspflicht nach der Schule, war der Kaiser für zwei wesentliche Entwicklungen verantwortlich, die vermutlich ohne die zentrale Weltmonarchie niemand möglich gewesen wären:

Erstens förderte von Monroe die Weltraumfahrt, die so viele Jahrzehnte brach gelegen hatte, da sie keinen kurzfristigen Profit abzuwerfen versprach. Mit den finanziellen Mitteln der ganzen Welt ausgestattet -- und diese waren und sind enorm! -- schuf er das neue Weltraumprogramm, das der wirtschaftlichen Erschließung der Himmelskörper des Sonnensystems dienen sollte. Und er versprach keine kurzfristigen Show-Effekte, wie damals die Mondlandung, sondern stellte das Ganze auf langfristige Beine und baute zuerst die notwendigen Basistechnologien auf, dann Weltraumstationen im Orbit um den Planeten und erst dann, nachdem das alles etabliert war, kamen die weiteren Ziele: Mond, Mars und Asteriodengürtel an die Reihe.

Zweitens rief von Monroe die \emph{Gesellschaft für Langfristforschung} ins Leben, eine rein aus Steuermitteln finanzierte Forschungseinrichtung die -- salopp ausgedrückt -- folgenden Auftrag hatte: ">Forscht in alle Richtungen, die a) extrem viel Geld kosten und b) \underline{keinen} kurzfristigen wirtschaftlichen Erfolg versprechen!"< Völlig zurecht, wie sich allerdings erst viel später herausstellte, war von Monroe davon ausgegangen, dass sich derartige Forschung irgendwann bezahlt machen würden. Ihm war klar, dass diese Gesellschaft Dinge entwickeln würde, auf die Niemand sonst kommen würde und dass sie ihm oder seinen Erben irgendwann Technologien in die Hand geben würde, die den Machterhalt sicherten.

Als er -- in sehr hohem Alter -- seine Macht an seinen Erben übergab, so stellte er sicher, dass -- neben den Grundlagen seiner Herrschaft -- auch diese beiden Institutionen die Zeit überdauern würden.




\chapter{Der erste Kontakt}

Captain Justin Svårregad saß in der Kanzel seines Aufklärers und las in einem Buch. Ihm war es augenblicklich völlig gleichgültig, dass er den Rang eines Captain hatte und als solcher natürlich eine Vorbildfunktion. Es gab ja schließlich niemand anderen an Bord seines Schiffes, für den er ein Vorbild sein könnte. Ja, die automatischen Aufzeichnungssysteme liefen und würden zeigen, dass er durchaus seine Job mit mehr Elan hätte ausführen können. Aber er hatte nicht vor, den Behörden nach der Landung einen Grund zu geben, die Daten auszuwerten. In Gesellschaft anderer war er stets tadellos, andernfalls hätte er es auch nicht zum Captain gebracht. Wenn er alleine war, gönnte er sich ein gewisses Maß an Freiheit.

Der kleine Maschine funktionierte wie am Schnürchen. Also bliebt ihm nichts anderes zu tun, als auf den nächsten Sprungpunkt zu warten -- und natürlich die Umgebung zu scannen, schließlich flog er einen Aufklärer. Zwar war er noch weit von seinem Operationsgebiet entfernt, aber sein Ausbilder in der kaiserlichen Akademie hatte ihm und seinen Kameraden immer wieder eingebläut: ">Meine Herren, ihr seid immer im Dienst, immer! Auch wenn ihr in der Passage seid, während Start und Landung, ja, auch in Eurer Freizeit! Das Operationsgebiet ist klein, der Weltraum ist groß. Merken Sie sich das!"< Bisher war er gut mit diesem Motto gefahren.

Also, routinemäßig und natürlich nicht allzu aufmerksam blickte er auf die Anzeigen der Sensoren, die den Raum um die Maschine abtasteten und erforschten. Bei aller weit entwickelter Technologie, die Fähigkeit des menschlichen Gehirns zur Mustererkennung war immer noch ungeschlagen auf seinem Gebiet. Die Computer würden Kandidaten für Auffälligkeiten melden, aber ein Mensch gar den letzten Ausschlag, ob etwas nur eine Störung oder eine handfeste Anomalie war.

In diesem Raumabschnitt war allerdings nichts zu erwarten. Er lag etwas Abseits der üblichen Routen, war aber nicht so unerforscht, als dass noch niemals jemand hier gewesen wäre.

Wie üblich meldeten die Systeme drei Kandidaten für ungewöhnliche Vorkommnisse. Es waren immer ca. zwei bis fünf, die gemeldet wurden und Svårregad argwöhnte, dass die Hersteller bei Motorola das genau so einstellten, um den Piloten bei Laune zu halten. Er blickte mit halben Auge zu den Anzeigen und wollte sie gerade quittieren, da stutzte er. Irgend etwas war nicht normal. Er besah sich die Daten und überprüfte den entsprechenden Raum genauer. Den zweiten \emph{Kandidaten} konnte er nicht so einfach als übliches Vorkommen im Raum abtun.

Der Raum war auf eine auffällige Weise gekrümmt. Die im Raum vorkommenden Gravitationsströme bogen, zerrten und knickten das Raum--Zeit--Kontinuum auf unterschiedlichste Arten und Weisen, aber das hier war seltsam. Er besah sich die Daten noch einmal und noch ein drittes Mal, dann löste er den Flottenalarm aus\dots

\mbox{}

Admiral Terschenkow blickte in die Runde. Gerade schloss sich die Tür zum letzten Mal, alle geladenen Offiziere waren anwesend.

">Meine Damen und Herren"<, begann er, ">ich mache es kurz. Heute Mittag, 10:76~terranischer Standard, wurde von einem Aufklärer, der im Sektor \textsf{27b} unterwegs war, der Flottenalarm ausgelöst. Sie werden gleich Gelegenheit bekommen, sich die Daten, die dem entsprechenden Captain aufgefallen waren, mit eigenen Augen zu besehen, aber vorher möchte ich Sie noch auf einen Umstand hinweisen: Obwohl sie alle mit dem Umgang geheimer Daten vertraut sind, so sind die Informationen, die Sie heute erhalten werden, in der höchsten Geheimhaltungsstufe violett eingestuft. Vielleicht muss ich dem einen oder der anderen noch einmal auf die Sprünge helfen, denn Stufe violett ist schließlich noch niemals zum Einsatz gekommen und Ihre Zeit seit der Akademie ist lange her"<. Er machte eine Pause. ">Außer Ihnen, mir, Captain Svårregad, dem Kaiser und einem Kommunikationsoffizier, der das Pech hatte, gerade im Dienst gewesen zu sein ist niemand, ich wiederhole \emph{niemand} über die Vorkommnisse informiert. Und das wird auch so bleiben, das garantiere ich Ihnen!"< Er blickte in die Runde.

Die Gesichter zeigten alle Schattierungen von Erstaunen, über Neugier bis hin zu Angst. Das war zu erwarten, die sechs anwesenden Offiziere waren zwar alle großartig und hatten sich vielfach bewährt, doch diese Art von Vorrede hatten sie noch nie gehört. Besser, jetzt etwas Angst zu empfinden, als später in Panik zu geraten. Panik war etwas, was sie sich absolut nicht leisten konnten, was Terschenkow aber bei keinem der Anwesenden erwartete.

Niemand stellte eine Frage. Sie warteten ab, was der Admiral ihnen eröffnen würde.

">Wie Sie wissen, erproben die terranischen Streitkräfte gerade neuartige Raumschiffe mit 5-D--Antrieb. Hierbei wird der Raum künstlich gefaltet und vorformt, um die Effizienz der Überlichtantriebe zu steigern. Die Resultate sind enorm vielversprechend. Sie wissen das, weil in ihren Einheiten erste Versuchsraumschiffe geflogen werden. Bitte, schauen Sie nicht so erstaunt, ich weiß, dass Sie alle annahmen, dass ausschließlich in Ihrer jeweiligen Einheit dieser Antrieb erprobt würde, doch kann die Gesellschaft nicht alles auf eine Karte setzen."<

Die \emph{Gesellschaft für Langfristforschung}, oder kurz \emph{GfL} hatte dieses neuartige Antriebselement entwickelt, nachdem sie schon vor über sechshundert Jahren den seit dem etablierten Überlichtantrieb entwickelt hatten, der die natürliche Raumkrümmung und die sogenannten Sprungpunkte ausnutzte und den Flug jenseits der Lichtmauer ermöglichte.

Einer der Anwesenden lächelte verstehend, Frau General Willencott von der Raumwaffe Nord schaute düster drein. Die anderen Gesichter blieben um ein gleichgültiges Aussehen bemüht.

Terschenkow fuhr fort: ">Sie wissen ebenfalls, dass diese künstlich verstärkten Raumkrümmung angemessen werden können und sich von natürlichen Phänomenen unterscheiden. Seit etwa 20 Jahren werden die Analyseeinheiten der Aufklärer mit Sensoren für derartige Signaturen ausgestattet, also lange bevor der Antrieb überhaupt einsetzbar wurde. Aber die \emph{GfL} hatte schon damals entsprechende Systeme fertiggestellt.

Sehen Sie selbst, was der Aufklärer \textsf{TCO-1960} aufgezeichnet hat!"<

Im Hologrammbereich des Konferenztisches erschien die Abbildung der Analyseeinheit des Aufklärers, nachgebildet und täuschend echt. Alle Anzeigen und Bildschirme waren detailliert zu erkennen, so dass jeder und jede der Anwesenden Gelegenheit hatte, eigene Schlüsse zu ziehen.

Schon nach Kurzen schnappte die erste nach Luft. Admiral Pamela Biko von der Heimatflotte war eine ausgezeichnete Analytikerin und hatte es als erste erkannt. Doch es dauerte nicht lange, da zeichnete sich auch auf allen anderen Gesichtern breites Erstaunen ab. Die Datenlage war klar, die Analyse einfach, die sich daraus ergebenden Schlüsse vielfältig und bedrohlich.

\mbox{}

%===============================
%===============================
{\sffamily\slshape\large TODO: Hier muss eine Beschreibung hin, wie Terschenkow die Besprechung plötzlich verlässt, um den Kaiser zu informieren, da Tutoro~II. gewillt ist, ihn zu empfangen. Er betritt den Kreml, den kaiserlichen Palast und denkt, dass es ziemlich doof ist, den Hof und die Regierung auf der Erde zu haben. Die Erde befindet sich seit der Expansion des Reiches eher in einem Randgebiet, Swoon wäre deutlich geeigneter. Tatsächlich hatte Swoon über 300 Jahre den Kaiserpalast beherrbergt, bis er vor etwa 100 Jahren wieder zurück zur Erde verlegt worden ist -- aus nostalgischen Gründen. Terschenkow hatte nichts für Nostalgie übrig, er war Pragmatiker. Dann bereitet er sich gedanklich auf das Treffen mit dem Kaiser vor, dass dann auch genauso abläuft, wie befürchtet\dots}
%===============================
%===============================

In obersten Raum des Palastes blickte Kaiser Tutoro~II. seinen Flottenadmiral an. ">Sind sie sicher?"< fragte er heiser.

Admiral Terschenkow nickte steif: ">Ja, Majestät. Alle ausgesuchten Offiziere des Krisenstabes sind zum gleichen Ergebnis gekommen. Natürlich wissen wir es nicht genau, aber wir müssen die Möglichkeit eines \emph{ersten Kontaktes} in Betracht ziehen, beziehungsweise, diese Möglichkeit an die erste Stelle setzen."<

">Das ist kaum vorstellbar"<, entfuhr es dem Kaiser. Admiral Terschenkow straffte sich innerlich. Er hasste es, wenn jemand die Realitäten nicht anerkannte, so irritierend sie auch sein mögen. Aber -- obwohl er durchaus Respekt für den amtierenden Kaiser empfand -- er konnte nicht umhin zu wünschen, dass der Großvater, Tutoro~I. Kaiser sein möge. Er war der Tradition gefolgt, dass der Kaiser der erste Soldat seines Volkes sein sollte. Sein Enkel vertrat eher die Meinung, dass der Kaiser der erste Bürger sein sollte, wobei \emph{Erster Bürger} insbesondere bedeutete, die meisten Privilegien zu besitzen.

Der Admiral schalt sich in Gedanken. Ganz so schlimm war Kaiser Tutoro~II. nun doch nicht. Natürlich genoss er eine herausragende Position und Vergünstigungen, aber es lebte nicht derartig in Saus und Braus, wie es andere Kaiser und Kaiserinnen von ihm getan hatten. Eigentlich war er ein guter Kaiser, ein guter oberster Verwalter des Reiches. Ob er sich allerdings in der derzeitigen Krise bewähren würde und das Reich anführen konnte, das bezweifelte Terschenkow.

Ächzend ließ sich Tutoro~II. auf einen Diwan fallen. ">Könnte es nicht doch noch andere Erklärungen geben?"< fragte er.

Terschenkow atmete leise durch: ">Die Daten, die Captain Svårregad uns übermittelt hat, zeigen deutlich und ohne Zweifel, dass die Raum--Zeit in den betreffenden Gebiet deutlich gekrümmt ist. Die Krümmungen entsprechen genau den Signaturen, die unsere neuen Antriebssysteme hinterlassen, die gerade von der \emph{GfL} entwickelt wurden und sich im Testeinsatz befinden. Sie stammen allerdings nicht von unseren Schiffen. Deren Positionen und Kursdaten kennen wir genau."<

">Also hat jemand anderes die gleiche Entdeckung gemacht."< folgerte der Kaiser.

">Das ist zwar denkbar, aber sehr unwahrscheinlich. Die \emph{GfL} verfügt über Forschungseinrichtungen und --budgets, die kein Privatmann besitzt. Hinter dieser Gesellschaft steht die gesamte Wirtschaftskraft des Reiches. Trotzdem hat es sehr lange gedauert, bis das Antriebssystem einsatzbereit gewesen ist und es hat unglaubliche Geldmengen verschlungen. Eine parallele Entwicklung können wir wohl schon alleine deshalb ausschließen."< Er machte eine Pause. ">Dazu kommt noch, dass die Signaturen etwa um den Faktor 10.000 größer sind, als die unserer Schiffe!"< Er lies die Worte wirken.

">Und das heißt?"< fragte der Kaiser. Die Wirkung war offenbar verpufft.

">Das heißt"<, fuhr Terschenkow fort und er wurde langsam ungeduldig, durfte sich das aber nicht anmerken lassen, ">dass die Antriebe, deren Signatur wir angemessen haben, um ziemlich genau diesen Anteil leistungsfähiger sind, als unsere."<

Der Kaiser lächelte und sagte: ">Sind Sie sicher, dass hier nicht verletzter Stolz im Spiel ist? Bitte bedenken Sie, ein Effizienzfaktor 10.000 von der ersten Idee, dem ersten Einsatz bis hin zur Perfektion ist nicht ungewöhnlich. Jemand hat uns überflügelt, jemand hat die \emph{GfL} überflügelt und Sie scheinen das nicht wahrhaben zu wollen. Gewiss, es ist unwahrscheinlich, aber doch nicht unmöglich. Es hat immer wieder zufällige Entdeckungen gegeben, die mit sehr wenig Aufwand das gleiche ermöglicht hatten, was in den großen Forschungszentren nur mit Mühe gelungen ist. Das alleine deutet für mich noch nicht auf einen \emph{Ersten Kontakt} hin."< Der Kaiser schien beruhigt zu sein.

Doch Terschenkow schüttelte nur bedächtig seinen Kopf: ">Leider ist es nicht ganz so, wie Sie vermuten, Euer Majestät. Der neuen Antrieb wurde von der \emph{GfL} zu einem Zeitpunkt der Flotte übergeben, wo die Effizienz schon sehr weit getrieben wurde. Schon mit dem heutigen Antrieb ist es möglich, die Galaxis in Gänze in ca.~30~Tagen zu durchqueren. Eine Effizienzsteigerung um den Faktor 10.000 bedeutet, dass der Antrieb in der Lage ist, den Leerraum zwischen den Galaxien zu durchqueren. Eine solche Effizienz benötigen Sie innerhalb der Galaxis nicht. Im Gegenteil: die Navigation ist bei solchen Geschwindigkeiten sehr kompliziert und langwierig, denn Sie brauchen auch sehr massive Sprungpunkte auf die Sie lange warten müssen. Nein, das ist keine \emph{normale} Effizienzsteigerung.

Sowohl nach meiner Meinung, als auch nach der Meinung der Expertenrunde glauben wir zu wissen, dass, wer immer in diesem Raumschiff sitzt, dass hier durch den Sektor \textsf{27b} geflogen ist, von außerhalb unserer Galaxis stammt. Wir haben es wirklich mit einem \emph{Outsider} zu tun, mit einer Lebensform von außerhalb unserer Milchstraße. Und das wäre definitiv ein \emph{Erster Kontakt}, nicht wahr, Majestät?"<

">Und es gibt keine andere Deutung der Datenlage?"< fragte Tutoro~II. noch einmal. Der Admiral seufzte innerlich. Dann sagte er mir fester Stimme: ">Keine, Majestät. Wir müssen uns auf eine Invasion außergalaktischer Intelligenzen einstellen."<

Der Kaiser richtete sich auf: ">Invasion? Wieso denn Invasion? Ich verstehe das mit den außergalaktischen Intelligenzen, auch wenn ich es immer noch sehr abenteuerlich finde. Aber wieso in aller Welt denken Sie an feindselige Aktivitäten?"<

Terschenkow wurde innerlich wieder ruhiger. Die erste Hürde war genommen, jetzt befand er sich auf sicherem Terrain. ">Nun ja, wir können das natürlich nicht mit Bestimmtheit sagen. Aber es gibt gute Gründe, von diesem Fall auszugehen.

Zum Einen hat zumindest die Geschichte unseres Volkes, der bislang einzigen intelligenten Lebensform, die wir kennen, gezeigt, dass der Vorstoß in neue Welten, neue Kontinente oder Sternensysteme meist das Ziel verfolgte, diese entdeckten Welten in Besitz zu nehmen, zu besiedeln oder wirtschaftlich auszubeuten. Eine evtl. vorhandene Bevölkerung wurde dabei meist verdrängt oder bekämpft. Wenn Sie die altterranische Geschichte studieren, so finden Sie dafür sehr viele Beispiele, von der Besiedelung der amerikanischen Kontinente oder der wirtschaftlichen Ausbeutung des indischen Subkontinents. Auch der Schritt in den Weltraum ist stets diesem Dogma gefolgt. Auf vielen Planeten, die einstmals eigene Lebensformen hervorgebracht haben, leben mittlerweile nur noch Lebensformen terranischen Ursprungs, von der Bakterie bis zum Homo sapiens."<

">Ich weiß das!"< unterbrach ihn der Kaiser, ">Aber wer sagt, dass das bei anderen Lebensformen genauso ist? Vielleicht sind sie ja friedlicher?"<

">Oh, auch die Besiedelung der Planeten, die heute das Reich formen, ist meist friedlich erfolgt. Schließlich sind wir nie auf intelligente Lebensformen gestoßen, die uns Widerstand leisten konnten. Aber es ändert nichts daran, dass wir heute die dominierende Rasse in der Galaxis sind. Als ich von Invasion sprach, so habe ich nicht definiert, dass sie bewusst und geplant aggressiv sein muss. Aber ich gehe lieber davon aus, denn: zweitens, es ist besser, vorsichtig zu sein, als später festzustellen, dass man das Ruder aus der Hand gegeben hat und\dots">

Der Kaiser unterbrach ihn erneut mit einer unwirschen Handbewegung: ">Also, das geht mir jetzt alles viel zu schnell und zu hoppla-hopp. Das Militär sieht überall Aggressoren, und Sie, mein verehrter Admiral, bilden da keine Ausnahme. Vielleicht fehlen Ihnen ja die Feinde, so dass Sie jetzt welche von Außen sehen, wo keine sind. Nein, das muss ich erst mit dem Kämmerer besprechen, dann sehen wir weiter."<

">Mylord, der Kämmerer ist für Geheimhaltungsstufe violett nicht freigegeben. Darf ich vorschlagen\dots"<

">Nein, Sie dürfen nicht! Wie Sie wissen, bin ich es immer noch persönlich, der die Geheimhaltungsstufe violett anordnet und freigibt. Und ich werde mich mit dem Kämmerer besprechen. In der Zwischenzeit\dots"<

Diesmal kam die Unterbrechung aus der Uniformtasche des Admirals. Das COM vibrierte deutlich sichtbar und gab leise, zwitschernde Töne wieder. Der Kaiser brach ab und blickte den Admiral fragen und sehr unzufrieden an. Admiral Terschenkow brach der Schweiß aus. Während einer kaiserlichen Audienz durfte ihn niemand stören, das wussten alle. Doch dann fiel ihm ein, dass er vor Beginn seinen COM so eingestellt hatte, dass nur noch Anrufe der Stufe violett erreichen konnte. Das bedeutete, jemand aus dem innersten Kreis hatte eine dringende Nachricht für ihn.

Er straffte sich und sagte: ">Majestät, das ist eine Dringlichkeitsbotschaft aus dem Krisenstab für mich. Sie duldet keinen Aufschub!"<

Der Kaiser seufzte und machte eine lässige Handbewegung: ">Aber bitte, macht schnell, und stellt laut, ich möchte hören, was so wichtig ist, dass ich gestört werden darf."< Es sprach der Monarch aus Tutoro~II.

Terschenkow griff in die Uniformjacke und holte das kleine, eiförmige Gerät heraus. Ein Druck auf die Seite und es war aktiviert: ">Ja, was ist so dringend?"< fragte er im Befehlston. Es antwortete der Kommunikationsoffizier, der schon die Übertragung der Captain Svårregad aufgefangen hatte und jetzt gezwungenermaßen zum Krisenstab gehörte: ">Admiral, wir haben jeden Kontakt mit Captain Svårregad verloren. Nicht nur, dass er sich nicht mehr meldet, auch die Telemetriedaten seines Aufklärers kommen nicht mehr an. Das Schiff ist spurlos verschwunden."<

">Was ist die letzte Telemetrie, die wir erhalten haben?"<

">Das Raumschiff beschleunigte unter Normalantrieb mit Höchstwerten. Es wirkt so, als hätte Svårregad versucht, zu entkommen, bevor das Schiff\dots"< er lies den Satz unvollendet.

Terschenkow blieb äußerlich ganz ruhig: ">Danke. Ich komme. Informieren Sie den Krisenstab. Treffen in 10 Minuten."< Dann schaltete er ab und stecken den COM wieder in die Jacke.

">Euer Majestät"<, sagte er ruhig, aber bestimmt, ">ich denke, hier haben wir ein Indiz für eine aggressive Handlung. Mit Eurer Erlaubnis werde ich alles für eine bevorstehende Invasion ausrichten. Höchste Geheimhaltungsstufe bleibt bestehen, keine Kommunikation zu niemandem. Darf ich mich entfernen?"<

Kaiser Tutoro~II. wirkte geschockt. Mit bleichem Gesicht sagte er, wie von weiter Ferne, ">Ja, natürlich. Ihr dürft gehen!"<

Und während Admiral Terschenkow dem Ausgang zuging wünschte er sich, nicht zum ersten oder letzten Mal, einen Kaiser mit etwas mehr militärischer Ausbildung oder mehr Initiative.

\chapter{Nachforschungen}

">Schätzchen, bist Du so weit?"< rief Frau deCartier aus dem Hinterzimmer.

">Ja"< rief Tamara zurück, ">Ich habe die Liste gefunden und bin bereit. Möchten Sie nicht doch mitkommen?"<

">Nein, danke"< antwortet Frau deCartier, ">Ohne mich bist Du doch schneller. Und heute brauche ich keine besondere Ablenkung, mein Bingo--Abend steht noch bevor."<

">Dann ist gut. Ich bringe ihnen die Sachen noch vorbei, dann bin ich weg. OK?"<

Frau deCartier streckte ihren Kopf durch die Tür: ">Ja, ist gut. Wir sehen uns dann Morgen. Ach nein, ich bin vergesslich, Morgen ist ja Samstag. Bis Montag dann also."<

">Ich schaue in jedem Fall noch einmal nach Ihnen, wenn ich mit den Einkäufen zurückkomme"<, versprach Tamara, dann verlies sie die Wohnung.

Tamara Svårregad absolvierte gerade ihr soziales Jahr. Nach Abschluss der Schule war für alle Kinder des Reiches ein soziales Jahr und ein Jahr Wehrdienst Pflicht. Die Reihenfolge konnte man wählen und Tamara hatte den Wehrdienst schon hinter sich. Die Zeit bei der Tiefenradaraufklärung war interessant gewesen, zumindest interessanter als die Grundausbildung. Aber auch das soziale Jahr machte ihr Spaß. Sie betreute alte Menschen, unterhielt sich mit ihnen, kaufte für sie ein und unternahm Botengänge. Das war zwar nicht notwendig, die Läden lieferten auch nach Hause, aber die alten Leute freuten sich sehr über den Besuch, denn häufig waren die Verwandten gar nicht mehr auf dem gleichen Planeten zu Hause und spätestens über 120~Jahren reiste man nicht mehr mit den Fernraumschiffen.

Frau deCartier war schon 140~Jahre alt, aber noch rüstig und interessant. Sie hatte viel zu erzählen aus ihrem langen Leben und Tamara hörte hier gerne zu. Heute aber freute sie sich schon sehr auf das Wochenende. Ihr Vater war zwar nicht im Lande, er musste einen mehrwöchigen Erkundungsflug machen, aber sie würde sich noch zu ihrer Mutter fahren und dann etwas Zeit mit ihren Freunden verbringen.

Schnell kaufte sie für Frau deCartier ein, lieferte die Sachen in der Wohnung ab und verabschiedete sich ins Wochenende, nachdem sie noch einmal nach der alten Dame geschaut hatte.

Zu Hause angekommen stellte sie sofort ihren \emph{Piepser} ein, ein kleines Gerät, das sich noch in der Schule im Elektronik--Kurs gebaut hatte. Es konnte Wellen auf unterschiedlichen Frequenzen des elektromagnetischen Spektrums oder des Überlichtbandes empfangen und in akustische Signale umsetzen. Im Astronomie--Kurs hatte sie das Gerät verwendet, um Quasare und Pulsare hörbar zu machen. Dann hatte sie Sonneneruptionen im elektromagnetischen und Überlichtband empfangen und aus dem Laufzeitunterschied die Entfernung Erde---Sonne berechnet.

Ihr Vater war sehr stolz gewesen. Als er dann immer wieder und auch für längere Zeit von zu Hause fort gewesen ist, hatte er ihr die Eigenfrequenz seines Raumschiffs genannt und sie hatten den Piepser darauf eingestellt. So hatte sie immer noch \emph{Kontakt} zu ihrem Vater, auch wenn er nicht da war. Mittlerweile war sie 20~Jahre alt und wohnte nicht mehr daheim. Doch aus Gewohnheit lies sie den Apparat immer noch mitlaufen, wenn sie wusste, ihr Vater war in den Tiefen des Weltraumes unterwegs.

Das kleine Gerät brauchte nach dem Einschalten immer ein paar Minuten, bis es sich komplett konfiguriert und die Überlichtantenne geladen hatte. Tamara ging derweil ins Schlafzimmer und zog sich um. Die eher konservative Kleidung, die sie für die Arbeit mit den alten Leuten trug, war nichts für die Bars und Clubs, in die sie heute noch zu gehen gedachte. Bevor sie noch in die Dusche ging, hörte sie die regelmäßigen Impulse, die aus dem Piepser drangen.

Unter den warmen, parfürmierten Wasserstrahlen der Dusche entspannte die junge Frau sich und wusch den Staub der Straßen von Sweet ab, der Hauptstadt von Swoon, dem ehemaligen Kaisersitz. Kurz noch durch den Trockner, dann konnte sie sich den schicken Rock anziehen, den sie letzte Woche entworfen hatte. Auf einmal stutzte sie. Etwas hatte sich verändert während sie sich erfrischt hatte. Dann fiel es ihr auf: es war still geworden. Aus dem Piepser drang kein Ton mehr.

\emph{Oh nein}, dachte sie, \emph{Jetzt ist das Ding kaputt. Wie blöd!} Jetzt würde sie doch wieder den ganzen Abend zu Hause bleiben, bis sie den Fehler gefunden hatte. Sie kannte sich. 

Tamara ging in den Flur, um nachzusehen. Die Kontrollleuchte an der Vorderseite leuchtete, ebenso, wie die Lampe der Feinabstimmung. Sie schaltete das Gerät auf Diagnose: alle Tests, die der kleine Apparat standardmäßig durchführte, um seine Funktionsweise zu testen, verliefen erfolgreich. Elektromagnetische und Subraumwellen konnten problemlos erfasst werden. Warum also zeigt es nichts an? Hatte der Piepser die Programmierung verloren?

Sie schaltete ihr COM ein, verband es mit dem Piepser und zeigte die aktuellen Einstellungen an. Das Gerät war völlig korrekt auf die Telemetrie des Aufklärers \textsf{TCO-1960} abgestimmt. Sie hätte weiterhin etwas hören müssen. Auch wenn der Flieger mittlerweile gelandet und abgestellt in der Wartungshalle gestanden hätte, hätte sie ein Signal reinbekommen. Doch das war nicht der Fall, sie wusste ja, dass ihr Vater unterwegs war. Aus dem Lautsprecher hätte weiterhin das beruhigende, regelmäßige Piepsen kommen müssen, klare Impulse von jeweils einer Zehntelsekunde Dauer, mit der die Maschine ihre kompletten Zustands--, Diagnose-- und Flugdaten übermittelte. Tamara konnte die Daten zwar mit dem Piepser nicht entschlüsseln, aber das gepulste Signal reichte ihr ja vollkommen.

Nun herrschte vollkommene Stille. Das hatte sie noch nie gehört. Es schien, als wäre die Flugmaschine ihres Vaters nie gebaut worden.

Tamara stand wie elektrisiert im Flur und regte sich nicht. Es konnte natürlich auch sein, dass die Maschine vollständig zerstört wäre. Auch das würde das  Ausbleiben der Telemetriedaten erklären. Aber war das möglich? Das Reich befand sich nicht im Kriegszustand, wer also sollte die Maschine abgeschossen haben? Vielleicht ein Unfall? Ja, das käme auch in Betracht, aber Tamara wusste, dass ihr Vater sehr umsichtig flog. Ein Unfall erschien ihr wenig wahrscheinlich.

Was also war geschehen? Sie musste es herausbekommen. Alle Pläne für das Wochenende waren vergessen, jetzt kam es darauf an, zu erfahren, was ihrem Vater und seiner \textsf{TCO-1960} widerfahren war. Eilig ging sie ins Nebenzimmer und setzte sich an die größere Konsole. Über dieses Gerät war sie, wie mit ihrem COM, in die zentralen Systeme des örtlichen Netzwerkes integriert und konnte Kontakt zum Reichsdatenverband aufnehmen, dem Galax--Net. Von hier führte ihr Weg sie in die astronomische Fakultät der kaiserlichen Akademie von Swoon. Hier hielt sie inne. Sie wusste ja überhaupt nicht, wo ihr Vater unterwegs gewesen war. Der Piepser konnte zwar die Richtung der Daten auswerten, die er empfing, tat das aber nur auf Anforderung. Tamara hatte bisher immer nur auf das regelmäßige Geräusch geachtet; es war ihr egal gewesen, wo ihr Vater flog, so lange er flog.

Erschüttert stellte sie fest, dass keine keine Möglichkeit und keine Idee hatte, etwas über den Verbleib ihres Vater herauszubekommen.

\mbox{}

Etwa zur gleichen Zeit, aber viele Lichtjahre entfernt auf der Erde traf sich der Krisenstab erneut. Wieder war es Admiral Terschenkow, der das Treffen eröffnete: ">Meine Damen, meine Herren. Der Kaiser hat sich unserer Empfehlung angeschlossen, die Angelegenheit als \emph{ersten Kontakt} mit möglicherweise invasorischen Absichten zu betrachten. Wir sollen und werden entsprechend verfahren."< Das entsprach zwar nicht ganz den Tatsachen, es würde aber die Operation erleichtern, wenn die anderen Mitgliedes des Stabes sich der Unterstützung Tutoros~II. sicher wähnten. Und streng genommen hatte der Kaiser das Vorgehen ja auch nicht abgelehnt.

Terschenkow fuhr fort: ">Es ist darüber hinaus meine Pflicht, Sie zu informieren, dass auch der Reichskämmerer nun in den Kreis der Geheimhaltungsstufe violett aufgenommen worden ist, auf besondere Anweisung seiner Majestät. Er ist aber nicht, ich wiederholt \emph{nicht} Mitgliedes dieses Krisenstabes, dessen Besetzung und Leitung mit obliegt. Habe ich mir klar ausgedrückt?">

Er blickte in sechsfaches Nicken.

">Nachdem das letzte Treffen durch meine Pflicht zur Unterrichtung seiner Majestät unterbrochen worden ist, kommen wir nun dazu, das weitere Vorgehen zu klären. Oberstes Ziel ist es, den Eindringling aufzubringen und seine Absichten zu erkunden, soweit das möglich ist. Wichtig ist vor allem die Frage, \emph{von wo} der Eindringling gekommen ist. Unter allen Umständen muss es verhindert werden, dass er Informationen über die hiesigen Verhältnisse erlangt und sich mit seiner Heimatbasis in Verbindung setzt. Es ist durchaus möglich, dass die Vernichtung des Eindringlings notwendig wird, sollten wir für uns feindliche Absichten erkennen."< Terschenkow blickte sie eindringlich an.

">Sie alle sind für diesen Stab ausgewählt worden, weil sie nicht nur hervorragende Offiziere sind, deren Ausbildung sie befähigen sollte, die aktuelle Krise zu meistern, sondern weil sie sich ebenfalls als ausgezeichnete Wissenschaftler bewährt haben. Es wird sicher in naher Zukunft notwendig sein, Spezialisten in die Gruppe Violett aufzunehmen, bis dahin sind Sie unsere Wissenschaftsteam. Halten Sie die Augen und Ohren auf und ergründen Sie die Natur der Bedrohung. Ich denke, dass es nicht notwendig ist, darauf hinzuweisen, dass das Reich noch niemals in seiner mittlerweile 1.000--jährigen Geschichte einer derartigen Krise ausgesetzt gewesen ist. Der Fortbestand des Reiches, wenn nicht sogar der menschlichen Kultur und Rasse ist in Gefahr. Das hier ist kein Spiel!"<

">Wissen wir"<, schnappte General Willencott zurück, ">Sie wiederholen sich, Admiral!"<

% Sloppypar-Umgebung, da die mbox sonst Gefahr läuft, über den Rand zu ragen
\begin{sloppypar}
">Dennoch kann der Umstand nicht oft genug betont werden. Fahren wir fort: Statusbericht jeden Mittag, 10:00~Uhr terranischer Standard Einsatzbesprechung, wenn möglich persönlich hier im Raum, ansonsten per COM, auf Kanal \mbox{\emph{\#LTviolett1960}}, Flottenserver drei, persönlicher Code! Vor Ihnen auf dem Tisch liegen Mappen, die Ihnen ihre jeweiligen Aufgabengebiete zuordnet. Bis heute Nachmittag 16:50 Uhr erhalte ich von Ihnen Vorschläge zur Analyse der Situation und des weiteren Vorgehens."<
\end{sloppypar}

Er machte ein kleine Pause: ">Meine Damen, meine Herren, es ist noch ein Umstand eingetreten, der die Sache verkompliziert: Der Aufklärer von Captain Svårregad, der uns die ersten Daten übermittelt hat, ist spurlos verschwunden, sein Telemetriesignal verstummte vor etwas 20 Minuten. Vorher fuhr Captain Svårregad noch die Triebwerke hoch, der Grund ist uns allerdings nicht bekannt. Admiral Biko!"<

Die angesprochene blickte auf.

">Admiral, Sie erhalten hiermit Order, den verlorenen Aufklärer wieder zu finden. Sie werden sich mit der Heimatflotte in das Raumgebiet \textsf{27b} begeben, genauer, zur letzten bekannten Position von \textsf{TCO-1960}. Dort werden Sie, Sie \emph{persönlich} die Messungen von Captian Svårregad wiederholen und akkuratere Ergebnisse liefern. Es muss verhindert werden, dass der Eindringling durch die Analyse des Aufklärers oder die Befragung des Captain zu viele Informationen über uns Menschen und da Reich erhält. Sehen Sie zu, dass Sie ihn finden. Wegtreten!"<

Alle Mitglieder des Stabes nahmen die Mappen auf und verließen den Raum. Terschenkow nickte Biko noch einmal zu, als sie den Raum verließ. Admiral Pamela Biko verließ schnellen Schrittes das Gebäude und rief per COM einen Gleiter, der sie umgehend zum Moskauer Flughafen brachte. Dort nahm sie den Flieger nach \emph{\textbf{!!!}}, wo ihre Heimatflotte stationiert war. Unterwegs gab sie dem Kapitän ihres Flaggschiffes den Befehl, die ganze Flotte klar zum Auslaufen zu machen und sich im  Raum oberhalb der Ekliptik zu sammeln. Sie käme persönlich zur \textsf{CONCORDIA} und würde die anlaufende Operation leiten.

Es war noch früh am Morgen, als sie in \emph{\textbf{!!!}} die Bodenstation der Ranke betrat. Direkt neben den zivilen Bahnen verliefen auf ein paar militärische Kabel, so viele geeignete Orte entlang des Äquators gab es schließlich nicht. Sie wurde bereits erwartet und betrat kurz nach Verlassen des Fliegers die Kabine der Kabelbahn. Ein freundlicher Leutnant der Marine half hier in die Liege. Statt den üblichen 2\,g, mit denen die zivilen Bahnen ihre Passagiere nach oben beförderten, verlangen die militärischen Bahnen fast 4\,g Andruck von ihren Mitfahrern. Daher war der Gebrauch von speziellen Andruckliegen notwendig. Das schränkte zwar die Bewegungsfreiheit ein, verkürzte aber den Weg in den Orbit auf weniger als ein Viertel.

Als die Bahn ein paar Stunden später die Raumstation erreichte, spürte Biko endlich wieder die Schwerelosigkeit. Sie hatte sie sehr vermisst. Etwas über ein halbes Jahr war sie nun auf der Erde gewesen und ihre Flotte hatte ohne sie im Leerraum operiert. Sie dankte dem Schicksal, das ihr wieder einen guten Grund gab, persönlich dabei zu sein, so bedrohlich die Situation auch sein mag.

Wenig später schwebte sie durch die geöffnete Luke der \textsf{CONCORDIA}. Hinter hier raste das Raumtaxi, das sie den letzten Weg gebracht hatte, zurück zur Station.

Der erste Maat begrüßte sie, als sie den Raumanzug öffnete: ">Willkommen an Bord, Admiral. Es ist gut, Sie wieder an Bord zu haben."<

">Danke, Leutnant. Ist alles bereit?"<

">Das Schiff ist klar zum Auslaufen. Die anderen Einheiten warten bereits wie vereinbart. Es wird Zeit für den Sprungpunkt."<

">Dann setzen Sie Segel und lichten den Anker"<, scherzte sie.

">Aye, aye"<, antwortete der erste Maat.

Lächelnd schwebten sie beide die Querachse zur Zentrale entlang. Kurz darauf beschleunigte das gewaltige Schiff und eilte dem Rest der Flotte entgegen. Die \textsf{CONCORDIA} erreichte ihre Position nur fünf Minuten vor dem Sprungpunkt. Das Timing von Admiral Biko war mal wieder perfekt gewesen.

\emph{So sollte es immer sein}, dachte sie, als das Schiff in den Subraum eintauchte und Sektor \textsf{27b} entgegenflog. 

\mbox{}

In Moskau neigte sich langsam der Nachmittag, es ging auf Abend zu. Der Kaiser kehrte von einer langweiligen, aber durchaus notwendigen Besprechung mit dem Ministerrat zurück in sein Quartier im Kreml. Ächzend ließ er sich auf einen Diwan fallen und streckte die Beine aus.

">Dieser verdammte Palupczok weiß überhaupt nicht, wie man ein Ressort führt"<, sprach er vor sich hin, ">es ist die totale Fehlbesetzung."<

">Mit Verlaub, dann ersetzt ihn doch"<, kam es von der Wandtür her. Der Kämmerer war eingetreten und trug ein Tablett mit der üblichen Auswahl von Getränken.

">Gerne, aber wer wäre denn besser geeignet? Diese Schwachkopf von Müller--Diepkau möchte dem Arbeitsministerium zu gerne die Landwirtschaft hinzufügen, aber damit kann er genau so wenig umgeben, wie Palupczok. Ich sähe gerne die Fanatubo als Leiterin des Ressorts, aber sie hat deutlich zu verstehen gegeben, dass sie keine Ressortleitung übernehmen möchte"<

">Das ist aber doch nicht ihre Entscheidung, oder, Majestät?"< fragte der Kämmerer und stellte das Tablett auf einem Tischchen neben dem Diwan ab.

">Naja, natürlich ist es ausschließlich meine Entscheidung, aber jemand, der eine Aufgabe nicht machen möchte, macht sie schlecht. Das Landwirtschaftsressort ist mit zu wichtig, als dass ich dort jemanden sehe, der es genauso in den Sand setzt, wie es Palupczok gerade tut."<

">Wenn es mir erlaubt ist\dots"<, sagte der Kämmerer.

">Natürlich, natürlich. Habt Ihr eine Idee?"< Der Kaiser griff nach einem fein geschliffenen Glas aus Rubin, in dem ein ebenso roter Wein schwappte.

">Nun"<, sagte der Kämmerer leise, ">Ich könnte mir vorstellen, dass es der Freiherr von Drohnenbrecht zu etwas bringen könnte."<

">Bitte, nicht immer der"<, sagte der Kaiser, ">Ihr schlagt ihn immer wieder vor. Nur weil er der Schwiegersohn Eurer Tochter ist, glaube ich nicht, er wäre geeignet."< Der Kaiser blickte seinen Kämmerer an: ">Es ist ein allzu durchsichtiger Versuch der Vetternwirtschaft, den ihr hier erneut versucht."< Er lächelte nachsichtig dabei.

">Euer Majestät sind zu gütig"<, sagte deCartier mit einer leichten Verbeugung. ">Wie wäre es des weiteren mit der Komtesse von Avernien? Sie hat ihren Planeten gut im Griff und der ist für die Versorgung des Reiches extrem wichtig. Erst neulich, erinnere ich mich, habt ihr die Komtesse explizit lobend erwähnt"<

">Das ist richtig. Ach, mein lieber Kämmerer, wenn ich Euch nicht hätte. Das ist ein guter Rat. Bitte, leitet alles in die Wege, dass dieser Nichtsnutz Palupczok noch Morgen seinen Platz räumt. Die Komtesse soll umgehend ihre Zelte auf Avernien abbrechen und zur Erde kommen. Sie gehört in den Ministerrat."< Der Kaiser sah für einen Moment erleichtert aus, dann sackte er wieder in sich zusammen. Kurz hielt er den Wein vor sich und betrachtete das Farbenspiel im Glas, dann leerte er das Getränk in einem Zug.

">Majestät"<, sagte der Kämmerer leise, der dieses Schauspiel wohl gesehen hatte, ">ich fürchte, das war nicht die einzige Sorge, die Euch plagt."<

">Ach, Kämmerer"<, seufzte Tutoro~II., ">Das ist nur zu richtig. Es gibt da noch etwas. Und es ist gut, dass ihr hier seid."< Er nahm seinen COM zur Hand und sprach hinein: ">Kaiserlicher Erlass: Kämmerer Iwan deCartier wird freigeschaltet für die Sicherheitsstufe Violett."<

">Was hat das zu bedeuten?"< fragte deCartier.

">Das bedeutet, dass ich es nicht einsehe, dass diese Terschenkow alles alleine macht. Er mag ja ein guter Admiral sein, aber ich denke nicht, dass die Angelegenheit alleine dem Militär obliegen darf!"<

">Was für eine Angelegenheit?"< fragte deCartier nun sehr neugierig geworden. Hier passierte etwas von oberstem Rang. Er wollte und musste mehr darüber erfahren.

\chapter*{Reststücke}

%Vielleicht später zu verwenden
Es war noch früh am Morgen, als sie in Washington auf dem Raumhafen ankam, doch die Sommersonne schickte schon ihre heißen Strahlen auf die Betonfläche des Armeeraumhafens. Admiral Biko verließ den Transportgleiter und eilte auf ihr Flaggschiff zu. Sie liebte es, ihr Schiff zu betrachten. Es war gewaltig groß, über zwei Kilometer lang und fast einen halben Kilometer hoch. Trotzdem sah es schlank und elegant aus, ein Walzenraumschiff der \textsf{B-}Serie, von manchen auch \emph{Bleistiftraumer} genannt. Seine stumpf--graue Außenhülle reflektierte kaum das Sonnenlicht, gegen die zivilen Schiffe der Linien- und Kreuzfahrtschiffe kam es optisch nicht an. Aber das war Pamela Biko egal, auf die glänzende, glitzernde Erscheinung eines dieser aufgedonnerten Turistenschlepper konnte sie gut verzichten.

Als sie die Rampe hinaufkam, wurde sie oben schon vom ersten Maat erwartet. Hinter hier donnerte die Luke zu und sie vernahm noch die Geräusche, mit denen der letzte Zugang raumtauglich gemacht wurde. 

%Überarbeiten!
Die Mannschaft war schon an Bord und so weit es ging informiert, der Rest der sog. \emph{Heimatflotte} war schon im Raum. Eilig schritt sie die Rampe hoch und noch während hinter ihr die Schotten zufuhren gab sie per COM dem Kapitän der \textsf{CONCORDIA} den Befehl zum Start.

Als sie im Kommandoraum ankam, lag die Erde schon hinter ihr und das Schiff strebte einem Punkt oberhalb der Ekliptik zu. Hier war die Flotte versammelt, man wartete auf das Flaggschiff und den Sprungpunkt. Bikos Timing war exzellent gewesen, die \textsf{CONCORDIA} traf nur fünf Minuten vor ein. Gemeinsam tauchte die gesamte Heimatflotte in den Subraum ein, der sie in etwa vier Stunden in den galaktischen Großraum \textsf{27b} bringen würde. Zeit genug, die Offiziere zu unterrichten und die Parameter für die Suche festzulegen.



\end{document}












